\documentclass[11pt]{article}
\usepackage[margin=0.9in]{geometry}
\usepackage{amsmath,amssymb,amsthm,booktabs,mathrsfs}
\usepackage[T1]{fontenc}
\usepackage{lmodern,microtype}
\usepackage[colorlinks=true,linkcolor=blue,urlcolor=blue]{hyperref}
\newtheorem{theorem}{Theorem}[section]
\newtheorem{lemma}[theorem]{Lemma}
\newtheorem{proposition}[theorem]{Proposition}
\newtheorem{corollary}[theorem]{Corollary}
\theoremstyle{remark}
\newtheorem{remark}[theorem]{Remark}
\newcommand{\norm}[1]{\left\|#1\right\|}
\title{Toward a four-dimensional Yang--Mills construction\\
\large Notebook 19: the physical background Gaussian in the source coordinates}
\author{Research working document}
\date{17 October 2026}
\begin{document}
\maketitle
\begin{abstract}
We construct the actual homogeneous current-step pivot chart near
the trivial field and a physical stationary background for the
next ordinary averaging constraint. The chart is local and leaves
the retained coordinates unchanged. Two minimum comparisons use
the previously constructed ordinary minimizer to identify the
current classical action on the next fiber. An explicit rooted
gauge sets the source's averaged contour variable to the identity.
The resulting reduced fluctuation Hessian is a small perturbation
of the flat retained precision, uniformly in earlier averaging
depth and periodic volume. Its formula keeps the nonlinear
constraint curvature and the conversion from relative exponential
to ordinary potential coordinates. This gives a uniform complex
covariance neighborhood for the physical small-background branch.
We also isolate the local measure and gauge factors needed by
the interacting integral.
\end{abstract}

\section{From a classical minimum to the source Gaussian}
\label{ps19-target}

The existence argument needs the Gaussian in the actual quantum
integration coordinates. The fine constrained inverse and the
Hessian of a minimized coarse action are useful inputs, but
neither alone identifies that Gaussian. The source also changes
coordinates, eliminates pivot bonds, and includes an exponential
gauge term. Away from the flat field these operations contribute
current-dependent second derivatives.

Here the current lattice has unit spacing, the previous fine
spacing is $M^{-1}$, and the next spacing is a fixed odd $L$.
The background is in a sufficiently small global potential ball.
The constants are uniform in $M=L^k$ and compatible periodic
volume, with the specified source block/interface geometry.
The result concerns physical stationary backgrounds and their
analytic extension near the trivial field. The original larger
orbit domain, freely prescribed independent currents, and mixed
large-field histories require further estimates.


% Notebook 19 module. No preamble. Prefix pc-.
% Uses norm, theorem/lemma/proposition/corollary and standard AMS packages.
\section{The physical current-step chart and Gaussian precision}
\label{pc-section}

The current lattice has spacing one, the original lattice spacing
is $M^{-1}$ with $M=L^k$, and the next spacing is fixed odd $L$.
We work on compatible periodic tori with the ordinary block and
corridor geometry. The original CMP109 choice has $L>11$; the
local arguments below work for odd $L\ge3$. Positive current bonds
and real color coordinates orthonormal for the normalized trace
are counted once. Current pairings use counting measure.
Transposes are bilinear after complexification; numerical-range
bounds use the Hermitian pairing.

Let $S_M(a)$ be the coupling-free Wilson action at the ordinary
homogeneous minimizer with current datum $V=e^{ia}$. Notebooks
15--16 give, on a uniform complex supremum ball,
\begin{equation}\label{pc-classical-input}
 m_M(a)=DS_M(a),\qquad
 \norm{m_M(a)}_\infty\le C\norm a_\infty,\qquad
 \norm{D^2S_M(a)-q_M}_\kappa\le C_\kappa\norm a_\infty.
\end{equation}
Here $q_M$ is the physical flat Maxwell precision. The norm
$\norm{\cdot}_\kappa$ is the maximum of the exponentially weighted
absolute row and column sums, with a fixed finite number of
color and orientation coordinates per cell. It is a matrix
algebra norm. We choose $\kappa$ below the common exponent of
Notebooks 07 and 16. Regrouping unit cells into fixed $L$-blocks
only changes fixed constants and the exponent.

Let $K_L$ insert the pivot variables in the flat rooted constraint,
and let $T_L$ be the averaged root-to-site contour derivative.
Notebook 17 proves
\begin{equation}\label{pc-flat-input}
 \begin{gathered}
 A_{M,0}=K_L^T(q_M+T_L^TT_L)K_L,\qquad
 \gamma_*I\le A_{M,0}\le B_*I,\\
 \norm{A_{M,0}}_\kappa+\norm{A_{M,0}^{-1}}_\kappa\le C_*.
 \end{gathered}
\end{equation}
Constants below may depend on $L$, dimension, the color
normalization and $\kappa$, but not on $M$ or periodic volume.

\subsection{The actual ordinary average and its relative coordinates}
\label{pc-source-maps}

The source formulas are CMP109, pp.~253--255 and 265--268,
(0.11)--(0.16), (2.3)--(2.13). Raw I, lines 46952--47000, and
Raw II, lines 9045--9095 and 73865--73945, spell out their
normalization. The following argument uses these explicit local
maps, not a general-background estimate quoted from the deposit.

Write $P_V(y,x)$ for the analytic group mean of the prescribed
shortest root-to-$x$ contour products in the block centered at
$y$. The mean is analytic near identity, its first derivative
there is the arithmetic mean, and it is equivariant at both
endpoints. With $V^u(b)=u(b_-)V(b)u(b_+)^{-1}$,
\begin{equation}\label{pc-contour-gauge}
 P_{V^u}(y,x)=u(y)P_V(y,x)u(x)^{-1}.
\end{equation}
All these contours are block-internal. For
$c=(y,y+Le_\mu)$, $x'=x+Le_\mu$, and the straight reference
transporter $R_c(V)=V([y,y+Le_\mu])$, the ordinary average is
\begin{equation}\label{pc-ordinary-map}
 \begin{split}
 U_x(V)&=P_V(y,x)V([x,x'])P_V(y+Le_\mu,x')^{-1},\\
 \mathcal M_L(V)(c)&=
 \exp\left\{L^{-d}\sum_{x\in B(y)}
      \log[U_x(V)R_c(V)^{-1}]\right\}R_c(V).
 \end{split}
\end{equation}
This is the ordinary rooted map, not the fixed-reference
corrected average. It reads two endpoint blocks and their
connecting strings, and obeys
\[
 \mathcal M_L(V^u)=
 \mathcal M_L(V)^{u|_{\text{block centers}}}.
\]

Define the source's left-relative coordinates
\begin{equation}\label{pc-relative-maps}
 \begin{split}
 F_V(B;c)&=\frac1i\log[
  \mathcal M_L(e^{iB}V)(c)\mathcal M_L(V)(c)^{-1}],\\
 G^{\rm rel}_V(B;y,x)&=\frac1i\log[
  P_{e^{iB}V}(y,x)P_V(y,x)^{-1}],\\
 A_V&=D_BF_V(0)=LQ_V^\flat,\qquad
 T_V=D_BG^{\rm rel}_V(0).
 \end{split}
\end{equation}
$A_V$ is a rectangular constraint derivative; the reduced
precision will be written $A_M(V)$. The source kernel convention is
\begin{equation}\label{pc-kernel-convention}
 (Q_V^\flat B)(c)=\sum_b L^{-d}Q_V^\flat(c,b)B(b).
\end{equation}

\begin{lemma}[One pivot per nonlinear equation]\label{pc-pivot-lemma}
Let $p(c)$ be the middle current bond on the line through the
centers of $c$, crossing their common interface. The full
nonlinear equation $F_V(B;c)$ reads $p(c)$ and no other pivot.
For $a=\log V/i$ sufficiently small,
\begin{equation}\label{pc-pivot-inverse}
 \pi_c(V)=Q_V^\flat(c,p(c))=I+O(\norm a_\infty),\qquad
 h_V(c)=L^{d-1}\pi_c(V)^{-1},\qquad A_Vh_V=I.
\end{equation}
Here $h_V$ inserts only at the corresponding pivots. The linear
insertion $K_V$ retains every nonpivot coordinate and sets
\begin{equation}\label{pc-linear-insertion}
 (K_Vz)(p(c))=-\pi_c(V)^{-1}
       \sum_{b\text{ nonpivot}}Q_V^\flat(c,b)z(b).
\end{equation}
Thus $A_VK_V=0$ and $K_1=K_L$.
\end{lemma}

\begin{proof}
No internal root contour contains a pivot. A string for another
coarse bond has a different orientation or a different pair of
endpoint blocks, and cannot contain $p(c)$. This proves the
nonlinear incidence assertion. At identity, precisely $L$
starting points in the straight-string average contain $p(c)$:
their transverse coordinates are fixed and their longitudinal
coordinate runs through the first block. Its coefficient in
$Q_L^\flat$ is therefore $L^{-d}$, and in $A_1$ it is
$L^{1-d}$. Analyticity gives the displayed inverse on a fixed
small neighborhood. The common factor $L^{1-d}$ cancels when
solving the linear equation, proving
\eqref{pc-linear-insertion}.
\end{proof}

\subsection{Exact straightening, finite range, and the cutoff}
\label{pc-chart-section}

Write $F_V(B)=A_VB+C_V(B)$, where $C_V$ starts at degree two.
The source straightening map is the solution of
\begin{equation}\label{pc-straightening}
 \widetilde D_V(B)=C_V(B-h_V\widetilde D_V(B)).
\end{equation}
For each output $c$, the right side depends on
$\widetilde D_V(c)$ alone, by Lemma~\ref{pc-pivot-lemma}.
A finite-dimensional contraction, with derivative zero at
$B=0$, solves this equation uniquely and analytically on a
uniform small ball. There are finitely many local cell types.
No iteration propagates support to another constraint equation.
Set
\begin{equation}\label{pc-exact-chart}
 \Phi_V(z)=K_Vz-h_V\widetilde D_V(K_Vz).
\end{equation}
Then
\begin{equation}\label{pc-chart-identities}
 F_V(\Phi_V(z))=0,\qquad
 \Phi_V(z)(b)=z(b)\quad(b\text{ nonpivot}),
\end{equation}
and this is the unique small solution for the specified retained
coordinates. It is therefore the actual source chart, rather
than a different parametrization of the same fiber. Its first
two derivatives are
\begin{equation}\label{pc-chart-derivatives}
 \begin{split}
 D\Phi_V(0)&=K_V,\\
 D^2\Phi_V(0)[u,v]
   &=-h_VD^2F_V(0)[K_Vu,K_Vv]
    =-2h_VC_V^{(2)}(K_Vu,K_Vv).
 \end{split}
\end{equation}
$C_V^{(2)}$ denotes the bilinear homogeneous coefficient whose
diagonal is the quadratic term, not the full Hessian.

For precision estimates it is useful to record how finite
dimensional analyticity gives uniform weighted bounds. Let
$N_L$ bound the number of scalar coordinates read by a local
output, $I_L$ the number of output neighborhoods containing a
fixed input coordinate, and $\ell_L$ their diameter. Cauchy on a
fixed smaller local polydisc bounds coordinate derivatives
through order three by $C_L$. In particular, for
$b=\norm a_\infty$,
\begin{equation}\label{pc-local-operators}
 \begin{split}
 \norm{K_V}_\kappa+\norm{h_V}_\kappa+\norm{T_V}_\kappa
      &\le C_{L,\kappa},\\
 \norm{K_V-K_L}_\kappa+\norm{h_V-h_1}_\kappa
       +\norm{T_V-T_L}_\kappa&\le C_{L,\kappa}b.
 \end{split}
\end{equation}
The first and second fluctuation derivatives of $\Phi_V$,
and its first and second background derivatives, have fixed
finite range and uniform coefficient sums on a smaller ball.
The corresponding second fluctuation derivative differs from
its identity-background value by $O(b)$, by integrating its
bounded background derivative. The same argument covers the
mixed derivatives of total order at most three.

More explicitly, if a local map $f_r$ has second coordinate
derivatives bounded by $C_L$, define the contracted matrix by
$\langle u,H_sv\rangle=\sum_rs_rD^2f_r[u,v]$. At most $I_L$
outputs contain a fixed input, and each has at most $N_L$
other inputs. Hence
\begin{equation}\label{pc-contracted-tensor}
 \norm{H_s}_\kappa\le
 C_LI_LN_Le^{2\kappa\ell_L}\norm s_\infty.
\end{equation}
This proves the required row and column bounds without a volume
factor. The first-derivative estimates use the same incidence
count.

After the source coupling scaling, the original relative link
field is $\Phi_V(gz)$. Exactly,
\begin{equation}\label{pc-cutoff-identity}
 \Phi_V(gz)(b)=gz(b)\quad(b\text{ nonpivot}).
\end{equation}
Thus a retained-link product of indicators $|B'(b)|<r$ becomes
$|z(b)|<r/|g|$, with no changed coordinate argument. CMP109,
p.~267, explains the resulting control of eliminated pivots;
CMP116, (2.2)--(2.3), explicitly uses the retained bond set.
Raw II's abbreviated description as a product over bonds must
not be used to discard other indicators. Any additional pivot,
contour or reconstructed-field characteristic function remains
at its exact pullback through $\Phi_V(gz)$. No equality between
such an enlarged window and a rectangular retained window is
claimed.

\subsection{The exponential gauge factor}
\label{pc-gauge-section}

With normalized trace $\tau(I)=1$, the analytic absolute gauge
function is
\begin{equation}\label{pc-gauge-function}
 \mathcal G(V)=\sum_y\sum_{\substack{x\in B(y)\\x\ne y}}
 \left\{1-\frac12\tau[P_V(y,x)+P_V(y,x)^{-1}]\right\}.
\end{equation}
On real unitary fields this is the source sum of
$1-\Re\tau P_V$. Its step factor is
$\exp[-g^{-2}\mathcal G(V)]$; field-independent normalization
and source contour indicators stay separate.

Suppose $P_V(y,x)=1$ for every pair. Since
$1-\tfrac12\tau(e^{iX}+e^{-iX})=\tfrac12\tau(X^2)+O(X^4)$,
\begin{equation}\label{pc-gauge-hessian}
 \begin{gathered}
 \mathcal G(V)=0,\qquad
 D_B\mathcal G(e^{iB}V)|_{B=0}=0,\\
 D_z^2\mathcal G(e^{i\Phi_V(z)}V)|_{z=0}
      =K_V^TT_V^TT_VK_V.
 \end{gathered}
\end{equation}
The nonlinear contour term does not contribute at quadratic
order because the first derivative of the penalty vanishes.
This is the averaged-contour penalty, not a tree penalty.

There is also an exact simplification. Contours are internal to
blocks and contain no pivots. If $E_f$ extends retained
coordinates by zero on pivots, then
\begin{equation}\label{pc-gauge-pivot-independence}
 \mathcal G(e^{i\Phi_V(z)}V)=\mathcal G(e^{iE_fz}V),\qquad
 T_Vh_V=0,\qquad T_VK_V=T_VE_f.
\end{equation}
The gauge factor is thus literally a local function of the
unchanged retained coordinates. Its cubic remainder after the
quadratic term is bounded by $C_L$ times the sum of cubed local
suprema of $|z|$. After $z\mapsto gz$ and division by $g^2$,
this is $O(g)$ times that local cubic sum, as long as
$|g|\norm z_\infty$ stays in the fixed chart radius.

At a general background with $P_V\ne1$, the absolute function
\eqref{pc-gauge-function} has a nonzero gradient and a
contour-curvature Hessian residual. These terms are retained in
its analytic extension and are $O(\norm a_\infty)$ perturbations
of the flat coefficients. The formal gradient times
$D^2\Phi_V$ vanishes by
\eqref{pc-gauge-pivot-independence}. Replacing the absolute
penalty by a relative one away from $P_V=1$ would change the
source function and is not done here.

\subsection{The physical stationary background}
\label{pc-stationary-section}

\begin{proposition}[Stationary homogeneous branch in the source gauge]
\label{pc-stationary}
For $W=e^{iw}$ with $\norm w_\infty$ sufficiently small, there
is an analytic current configuration $V_*(W)$ satisfying
\begin{equation}\label{pc-stationary-identities}
 \mathcal M_L(V_*)=W,\qquad P_{V_*}(y,x)=1,\qquad
 \norm{\log V_*/i}_\infty\le C_L\norm w_\infty.
\end{equation}
For real data it is a constrained local minimum of the
classical action, and a stationary point of the action plus
the absolute exponential-gauge function.
\end{proposition}

\begin{proof}
Apply Notebook15 with final spacing $L$, original spacing
$1/M$ and total ratio $ML$. Under $x'=x/L$, this is the
already established unit-final-spacing theorem. In dimension
four the coupling-free Wilson action has no extra scale
factor; physical potential and derivative norms change only
by fixed powers of $L$. Let $U_f(W)$ be that ordinary
minimizer and set
\[
 V_p(W)=\mathcal M^k(U_f(W)).
\]
Exact composition gives $\mathcal M_L(V_p)=W$.
Notebook14's ordinary trajectory and Notebook15's potential
bound give $\norm{\log V_p/i}_\infty\le C_L\norm w_\infty$
and a fine physical $N_\alpha$ bound $C_L\norm w_\infty$.

Let $U_c(V)$ be the current-scale ordinary minimizer. Choose a
common current comparison radius $\delta$ and then the data
radius so that $U_f(W)$ is in the current minimum ball,
every $U_c(V)$ with $\norm{\log V/i}_\infty<\delta$ is in
the final-spacing-$L$ minimum ball, and
$\norm{\log V_p/i}_\infty<\delta/2$. These are finitely many
conditions $C_L\delta<r_0$ and
$C_L\norm w_\infty<\delta/2$, independent of $M$ and volume.
Current minimality gives
\[
 S_M(\log V_p/i)\le\mathcal S(U_f(W)).
\]
The reverse inequality follows because $U_c(V_p)$ is a full
feasible competitor by ordinary composition and lies in its
common minimum ball. Hence equality holds. For any $V$ in the
current comparison ball with $\mathcal M_L(V)=W$, the same
full-minimum comparison gives
\begin{equation}\label{pc-minimum-comparison}
 S_M(\log V/i)\ge\mathcal S(U_f(W))=S_M(\log V_p/i).
\end{equation}
This proves the required constrained minimum directly.

Define a rooted current-site gauge by
$u(y)=1$ at centers and $u(x)=P_{V_p}(y,x)$ in $B(y)$.
It is local analytic and
$\norm{\log u/i}_\infty\le C_L\norm{\log V_p/i}_\infty$.
Set $V_*=V_p^u$. Equations
\eqref{pc-contour-gauge} and \eqref{pc-ordinary-map} give
$P_{V_*}=1$ and $\mathcal M_L(V_*)=W$ exactly. The local
coarse gauge invariance of Notebook16, with its smooth fine
extension and stricter common margins, gives
$S_M(\log V_*/i)=S_M(\log V_p/i)$. Therefore
\eqref{pc-minimum-comparison} applies with $V_*$ as well.
The gauge function has zero gradient there by
\eqref{pc-gauge-hessian}.

All the maps are analytic on a uniform complex ball. The
constraints and stationarity identities extend from the real
form by holomorphic uniqueness. No complex minimum
interpretation is asserted.
\end{proof}

\subsection{Both chart-curvature contributions to the Hessian}
\label{pc-hessian-section}

Let $a=\log V/i$ and define the ordinary potential in the
source relative chart by
\begin{equation}\label{pc-ordinary-coordinates}
 \Xi_a(B)=\frac1i\log(e^{iB}e^{ia}),\qquad
 a_V(z)=\Xi_a(\Phi_V(z)).
\end{equation}
For the bondwise left Jacobian
$J_aX=\int_0^1\operatorname{Ad}(e^{ita})X\,dt$,
\begin{equation}\label{pc-log-jacobian}
 D\Xi_a(0)=J_a^{-1},\quad
 J_a^{-1}=I+O(\norm a_\infty),\quad D^2\Xi_a(0)=O(\norm a_\infty).
\end{equation}
The last bound follows from analyticity and $\Xi_0(B)=B$.
Put $K_a=J_a^{-1}K_V$. Then
\begin{equation}\label{pc-second-coordinate}
 \begin{split}
 Da_V(0)&=K_a,\\
 D^2a_V(0)[u,v]&=J_a^{-1}D^2\Phi_V(0)[u,v]
                 +D^2\Xi_a(0)[K_Vu,K_Vv].
 \end{split}
\end{equation}
Thus, at a gauge-zero physical background $V_*(W)$, the exact
retained-coordinate Hessian of action plus gauge is
\begin{equation}\label{pc-physical-precision}
 \begin{split}
 A_M(V)&=K_a^TH_M(a)K_a+\operatorname{Curv}_M(V)
                  +K_V^TT_V^TT_VK_V,\\
 \langle u,\operatorname{Curv}_M(V)v\rangle
      &=\langle m_M(a),D^2a_V(0)[u,v]\rangle,\qquad
 H_M(a)=D^2S_M(a).
 \end{split}
\end{equation}
The two terms in \eqref{pc-second-coordinate} are respectively
the nonlinear constraint curvature and the conversion from
relative exponential to ordinary log coordinates. Both occur
in \eqref{pc-physical-precision}.

For an explicit check against the source current vertex, set
\[
 j_{\rm rel}=J_a^{-T}m_M(a),\qquad
 H_{\rm left}=J_a^{-T}H_M(a)J_a^{-1}
                     +\langle m_M(a),D^2\Xi_a(0)\rangle.
\]
Then the same precision is
\begin{equation}\label{pc-source-current-form}
 \begin{split}
 \langle u,A_M(V)v\rangle={}&
 \langle K_Vu,H_{\rm left}K_Vv\rangle
 +\langle T_VK_Vu,T_VK_Vv\rangle\\
 &-2\langle j_{\rm rel},h_VC_V^{(2)}(K_Vu,K_Vv)\rangle.
 \end{split}
\end{equation}
This has the sign and factor of Raw I's source vertex.
Differentiating the same action through a fine minimizing
extension writes the physical current as $H_1^TJ$; pure gauge
pieces pair to zero by gauge invariance. No independent $J$
is introduced here. Constrained stationarity gives
\begin{equation}\label{pc-current-stationarity}
 K_V^Tj_{\rm rel}=0,\qquad
 j_{\rm rel}=A_V^T\lambda,\quad \lambda=h_V^Tj_{\rm rel}.
\end{equation}
The second identity follows because $I-h_VA_V$ has range
$\ker A_V$. It does not make $j_{\rm rel}$ zero and does
not eliminate its contraction with $h_VC_V^{(2)}$.

\subsection{Uniform precision, covariance, and uniqueness bounds}
\label{pc-bounds-section}

\begin{theorem}[Physical-background retained precision]
\label{pc-precision-theorem}
For the background of Proposition~\ref{pc-stationary}, with
$b=\norm{\log V/i}_\infty$ sufficiently small,
\begin{equation}\label{pc-precision-perturbation}
 \norm{A_M(V)-A_{M,0}}_\kappa\le C_{L,\kappa}b.
\end{equation}
The full Hessian defined using the actual absolute gauge
function has the same analytic estimate on a full small
current-background ball. At arbitrary backgrounds it need
not be centered, since its first derivative can be nonzero.
For
\begin{equation}\label{pc-radius}
 C_{L,\kappa}b\le\gamma_*/2,\qquad
 C_*C_{L,\kappa}b\le1/2,
\end{equation}
one has
\begin{equation}\label{pc-inverse-bounds}
 \begin{gathered}
 \operatorname{Re}\langle z,A_M(V)z\rangle_{\rm Herm}
       \ge\tfrac12\gamma_*\norm z_2^2,\qquad
 \norm{A_M(V)^{-1}}_\kappa\le2C_*,\\
 \norm{A_M(V)^{-1}-A_{M,0}^{-1}}_\kappa
       \le2C_*^2C_{L,\kappa}b.
 \end{gathered}
\end{equation}
The constants are uniform in all preceding classical depths
$M=L^k$ and compatible periodic volumes.
\end{theorem}

\begin{proof}
Equations \eqref{pc-local-operators} and
\eqref{pc-log-jacobian} give
$\norm{K_a-K_L}_\kappa\le C_{L,\kappa}b$, with all
unperturbed factors uniformly bounded. Apply the weighted
algebra inequality and \eqref{pc-classical-input} to the
first term in \eqref{pc-physical-precision}. Its difference
from $K_L^Tq_MK_L$ is $O(b)$. Apply
\eqref{pc-contracted-tensor} to
\eqref{pc-second-coordinate}, using
$\norm{m_M(a)}_\infty\le Cb$, to obtain
$\norm{\operatorname{Curv}_M(V)}_\kappa\le C_{L,\kappa}b$.
The ordinary-log conversion part is even $O(b^2)$, but that
improvement is unnecessary. The gauge term differs from its
flat value by $O(b)$ through \eqref{pc-local-operators}.
The absolute gauge residual off $P_V=1$ is a local analytic
$O(b)$ term, so the same bound holds for that extension.

The Schur bound controls the Hilbert operator norm. Together
with \eqref{pc-flat-input}, it proves the accretivity assertion.
The other bounds follow by the Neumann series in the weighted
matrix algebra.
\end{proof}

For clarity, uniqueness of the source critical configuration
can be obtained on a uniform common coordinate neighborhood,
not merely from the Hessian at a point. Keep one fixed chart
$\Phi_V$ and let $z$ vary. Local Cauchy gives
$D\Phi_V(z)=K_L+O(b+\norm z_\infty)$, uniformly bounded
second derivatives, and
$\norm{a_V(z)}_\infty\le C_L(b+\norm z_\infty)$.
The same chain-rule proof, including the absolute gauge
contour residual, gives
\begin{equation}\label{pc-uniform-convexity}
 \left\|D_z^2\{S_M(a_V(z))+
       \mathcal G(e^{i\Phi_V(z)}V)\}-A_{M,0}\right\|_\kappa
       \le C_{L,\kappa}(b+\norm z_\infty).
\end{equation}
On a smaller real convex supremum ball this Hessian is bounded
below by $\gamma_*/2$. There is at most one critical point in
that ball. Every sufficiently small configuration on the same
fiber belongs to this chart, by uniqueness for its retained
relative coordinates. Thus the critical orbit and gauge-zero
selection in CMP109 agree with the constructed branch whenever
both are in this common homogeneous neighborhood. No global
minimum outside that neighborhood is inferred.

For any specified $\varepsilon>0$, the additional restriction
$b\le\varepsilon\gamma_*/C_{L,\kappa}$ gives
\begin{equation}\label{pc-relative-precision}
 \norm{A_{M,0}^{-1/2}[A_M(V)-A_{M,0}]A_{M,0}^{-1/2}}_{2\to2}
       \le\varepsilon.
\end{equation}
The counting covariance in the current retained coordinates
is $A_M(V)^{-1}$; no factor $M^4$ is inserted.

\begin{corollary}[The tunable covariance neighborhood]
\label{pc-tunable-covariance}
Let $C_0=A_{M,0}^{-1}$, so its real spectral bounds are
$\lambda_{\rm cov}=B_*^{-1}$ and
$\Lambda_{\rm cov}=\gamma_*^{-1}$. Fix any
$0<\theta\le1/4$ before the coupling is chosen. If, in addition
to \eqref{pc-radius}, the background radius satisfies
\begin{equation}\label{pc-covariance-radius}
 2C_*^2C_{L,\kappa}b\le\frac{\theta}{2B_*},
\end{equation}
then
\begin{equation}\label{pc-covariance-close}
 \norm{A_M(V)^{-1}-C_0}_{2\to2}
       \le\theta\lambda_{\rm cov}/2.
\end{equation}
This is exactly the covariance-neighborhood hypothesis of
Notebook18's tunable marked theorem. A common covariance
majorant and exponentially weighted row sum are uniform in
$M$ and volume.
\end{corollary}
\begin{proof}
Use \eqref{pc-inverse-bounds} and the Schur operator bound.
The same weighted inverse estimate gives the entry majorant
$J_{ij}=2C_*e^{-\kappa d(i,j)}$. For every smaller exponent
$0<\kappa'<\kappa$, its weighted row sum is uniformly finite
by the polynomial growth of the lattice, including periodic
distance on compatible tori. The real baseline has the stated
spectral bounds by \eqref{pc-flat-input}.
\end{proof}
The local interaction hypotheses of that marked theorem are
still separate. Its strict-background-locality version further
requires a support-valued precision expansion and an appropriate
product background domain, which the present kernel estimate
does not imply. Alternatively,
\eqref{pc-relative-precision} gives relative covariance error
at most $\varepsilon/(1-\varepsilon)$; the direct condition
\eqref{pc-covariance-radius} matches the stated theorem without
that conversion.

At the centered physical background, the exact classical
exponent after scaling is
\begin{equation}\label{pc-scaled-exponent}
 \begin{split}
 &-g^{-2}\{S_M(a_V(gz))+\mathcal G(e^{i\Phi_V(gz)}V)\}\\
 &\hspace{1cm}=-g^{-2}S_M(a)-\tfrac12\langle z,A_M(V)z\rangle
                  +\text{terms of positive order in }g,
 \end{split}
\end{equation}
locally for $|g|\norm z_\infty$ below the chart radius. The
linear term vanishes by physical constrained stationarity and
the gauge-zero condition. Haar density, the straightening
Jacobian and the previously integrated action correction
remain in the interaction; a coupling-suppressed quadratic
Haar term is not asserted to vanish.

\subsection{The remaining source interfaces}
\label{pc-scope}

The chart, gauge penalty, cutoff-coordinate identity and
precision comparison above concern the exact homogeneous
ordinary branch already matched in Notebooks15--17. They
do not extend the source's independent-free-current family:
the current in \eqref{pc-source-current-form} is the physical
action derivative. Nor do they establish the full original
orbit and polymer-restriction hypotheses of CMP116, or a
mixed-history boundary theorem.

Finite range of $\Phi_V,h_V,K_V,T_V$ means finite dependence
on the current background $V$. The minimizing response
$W\mapsto V_*(W)$ need not have strict finite support.
Weighted fluctuation-kernel locality in
\eqref{pc-precision-perturbation} is not a strict
support-valued expansion in $W$ or in every original fine
link. A complete localized expansion of the classical
remainder and the Haar, Jacobian and prior-action terms is
still a separate interaction theorem. Additional source
cutoff factors, when present, retain their exact pullbacks.

% No preamble. Local measure factors for Notebook19. Prefix pm-.
\section{Local Haar, straightening, and exponential-gauge interactions}
\label{pm-section}

Use the exact current-step chart of Section~\ref{pc-chart-section}.
The next block side $L$ is fixed, the current spacing is one,
and constants are uniform in the preceding homogeneous depth $M$
and compatible periodic volume. Backgrounds $V=e^{ia}$ range
over a sufficiently small product ball in the independent
current link coordinates $a(b)$. A slightly larger fixed ball
is used for Cauchy estimates. This section localizes the measure
factors; the classical and previously integrated action
remainders are separate.

Write $A_{c,p}=D_{p(c)}F_V(0;c)=L^{1-d}\pi_c(V)$ and define
the square ambient straightening map
\begin{equation}\label{pm-ambient-map}
 \psi_V(B)=B-h_V\widetilde D_V(B),\qquad
 F_V(\psi_V(B))=A_VB,\qquad
 \Phi_V(z)=\psi_V(K_Vz).
\end{equation}
Its retained components equal those of $B$. The nonlinear
equation for its pivot $p(c)$ reads that pivot and no other
pivot. All maps and derivatives used below are local analytic
functions on fixed neighborhoods.

More explicitly, choose supports $Y_r$ of current cells for
every primitive term below. Include both its background and
retained fluctuation dependencies. There are fixed constants
$s_L,r_L,n_L$ such that
\begin{equation}\label{pm-support-bounds}
 |Y_r|\le s_L,\qquad \operatorname{diam}Y_r\le r_L,
 \qquad \ell(Y_r)\le s_Lr_L,\qquad
 \sup_x\#\{r:x\in Y_r\}\le n_L.
\end{equation}
Here $\ell$ is shortest-tree length; periodic multiplicities
are retained. Finite enlargement of the two-block chart
neighborhoods and the block-internal contour neighborhoods
gives these constants. The local implicit equations do not
propagate support between different coarse equations.

For nonnegative $a_0,\kappa$ define the primitive local norm
\begin{equation}\label{pm-local-norm}
 \|l\|_{a_0,\kappa,R}
 =\sup_x\sum_{r:x\in Y_r}
 e^{a_0|Y_r|+\kappa\ell(Y_r)}
 \sup_{V\in\mathcal D_b,\,\|z|_{Y_r}\|_\infty\le R}|l_r(V,z)|.
\end{equation}
The fluctuation supremum may be taken over complex coordinates,
so it also bounds the fixed real cutoff. In particular,
\begin{equation}\label{pm-incidence-bound}
 \|l\|_{a_0,\kappa,R}
 \le n_Le^{a_0s_L+\kappa s_Lr_L}\sup_r\|l_r\|_\infty.
\end{equation}
Adding primitive terms with the same support preserves this
bound by the triangle inequality.

\subsection{The exact ambient Jacobian and the delta normalization}
\label{pm-jacobian-section}

\begin{lemma}[Straightening determinant]\label{pm-jacobian-lemma}
On a fixed small complex chart,
\begin{equation}\label{pm-jacobian-identity}
 \det D\psi_V(B)=\prod_c
 \frac{\det A_{c,p}}
      {\det D_{p(c)}F_V(\psi_V(B);c)}.
\end{equation}
Each ratio is analytic, equals one at $B=0$, and has fixed
local dependence.
\end{lemma}
\begin{proof}
Order the variables as retained coordinates followed by
pivots. The retained rows of $D\psi$ are $[I,0]$.
The pivot-to-pivot block is block diagonal, because each local
implicit equation reads only its own input pivot. Differentiating
$F_V(\psi_V(B))=A_VB$ in that pivot gives
\[
 D_{p(c)}F_V(\psi_V(B);c)
 D_{B(p(c))}\psi_V(B)(p(c))=A_{c,p}.
\]
Taking the determinant of this block triangular matrix proves
\eqref{pm-jacobian-identity}. At zero, $DF_V=A_V$ and
$D\psi=I$. All local blocks stay invertible on a smaller
fixed domain.
\end{proof}

Let $\operatorname{Tr}_{\mathfrak g}$ denote the ordinary
matrix trace in color-coordinate space. The logarithm of the
normalized $c$ factor is
\begin{equation}\label{pm-jacobian-log}
 j_c^{\rm str}(V;B)=
 -\operatorname{Tr}_{\mathfrak g}\log\left[
 I+A_{c,p}^{-1}
 \{D_{p(c)}F_V(\psi_V(B);c)-A_{c,p}\}\right],
\end{equation}
with branch zero at $B=0$. This is a determinant identity:
the diagonal Jacobian block itself is $(DF_p)^{-1}A_p$,
and no commutativity of color matrices is asserted.
Local analyticity and the logarithm series imply, for a fixed
$\rho_L>0$ and $|g|R<\rho_L$,
\begin{equation}\label{pm-jacobian-bound}
 \|j^{\rm str}(V;K_Vgz)\|_{a_0,\kappa,R}
 \le C_{L,a_0,\kappa}|g|R.
\end{equation}
Indeed, the matrix perturbation inside the logarithm is at
most $C_L|g|R$, which may be made at most $1/2$.

The determinant in Lemma~\ref{pm-jacobian-lemma} belongs to
the square ambient change of variables, not to the rectangular
derivative $D\Phi_V$. To identify its measure role, for real
backgrounds and test functions supported within the chart,
\begin{equation}\label{pm-delta-reduction}
 \begin{split}
 \int dX\,\delta(F_V(X))f(X)
 &=\mathcal D(V)\int dz\,
 \exp\left\{\sum_cj_c^{\rm str}(V;K_Vz)\right\}
 f(\Phi_V(z)),\\
 \mathcal D(V)&=\prod_c(\det A_{c,p})^{-1}.
 \end{split}
\end{equation}
First set $X=\psi_V(B)$ and then integrate the linear
constraint $A_VB=0$ over pivots. The pivot determinants have
the positive orientation of $L^{1-d}I$ on this real
neighborhood. A group delta at identity contributes only its
fixed exponential-coordinate normalization at $F=0$.
After retained variables are scaled by positive real $g$,
the additional factor $g^{\dim z}$ is field independent.
Complex-background formulas are holomorphic continuations
of the local densities; no complex distributional delta is used.

The factor $\mathcal D(V)$ is independent of fluctuations
but can depend on the prescribed next-scale field through $V$.
It is kept outside the fluctuation interaction, not discarded.
Its normalized logarithm is local:
\begin{equation}\label{pm-background-determinant}
 \log\frac{\mathcal D(V)}{\mathcal D(1)}
 =-\sum_c\operatorname{Tr}_{\mathfrak g}
       \log(A_{c,p}A_{1,c,p}^{-1}).
\end{equation}
Each term is $O(\|a\|_\infty)$ and has no small coupling
factor $g$.

\subsection{The relative Haar density}
\label{pm-haar-section}

Let $j(X)$ be the pullback density of Haar measure under
$X\mapsto e^{iX}$, divided by its value at zero. Near zero
it is nonvanishing and real analytic, with $j(0)=1$.
Haar invariance under inversion gives $j(-X)=j(X)$ on
the real form and hence by analytic continuation. Its local
logarithm therefore satisfies
\begin{equation}\label{pm-haar-even}
 \log j(0)=D\log j(0)=0,\qquad |\log j(X)|\le C|X|^2.
\end{equation}
This uses the positive real density and then its analytic
branch; it does not complexify an absolute determinant.
Translation by $V$ preserves Haar measure, so the relative
parametrization $e^{iX}V$ has the same density.

The exact Haar logarithms after scaling are
\begin{equation}\label{pm-haar-log}
 j_b^{\rm H}(V;gz)=\log j(\Phi_V(gz)(b)).
\end{equation}
On retained bonds the argument is exactly $gz(b)$; on pivots
it is the local chart function. From
\eqref{pm-haar-even}, the uniform local chart bound, and
\eqref{pm-incidence-bound},
\begin{equation}\label{pm-haar-bound}
 \|j^{\rm H}(V;gz)\|_{a_0,\kappa,R}
 \le C_{L,a_0,\kappa}|g|^2R^2.
\end{equation}
Any unnormalized value $j(0)$ is field independent and remains
in the overall normalization. A generally nonzero quadratic
Haar term remains in the interaction with coefficient $g^2$.

\subsection{The exact retained-coordinate gauge remainder}
\label{pm-gauge-section}

For a root-to-site pair $r=(y,x)$ define
\begin{equation}\label{pm-gauge-primitive}
 f_r(V;z)=1-\tfrac12\tau\left[
 P_{e^{i\Phi_V(z)}V}(y,x)
 +P_{e^{i\Phi_V(z)}V}(y,x)^{-1}\right].
\end{equation}
The source contours are internal to blocks and contain no
pivots. By \eqref{pc-gauge-pivot-independence}, this is
exactly the same expression with $\Phi_V(z)$ replaced by
$E_fz$, where $E_f$ inserts zero at pivots. Consequently
it is a local function of unchanged retained coordinates.

On the full small product background domain define
\begin{equation}\label{pm-gauge-remainder}
 \begin{split}
 v_r^{\rm G}(V;g,z)=-g^{-2}\bigl[
 &f_r(V;gz)-f_r(V;0)-gD_zf_r(V;0)z\\
 &-\tfrac12g^2D_z^2f_r(V;0)[z,z]\bigr].
 \end{split}
\end{equation}
The singularity at $g=0$ is removable, with value zero.
The finite-dimensional Taylor remainder gives
\begin{equation}\label{pm-gauge-bound}
 \|v^{\rm G}(V;g,z)\|_{a_0,\kappa,R}
 \le C_{L,a_0,\kappa}|g|R^3,
 \qquad |g|R<\rho_L.
\end{equation}
At a physical background $P_V=1$, both $f_r(V;0)$ and
$D_zf_r(V;0)$ vanish. Thus this is precisely the logarithmic
gauge interaction after its quadratic contribution has been
included in the precision. Elsewhere it is the exact
third-and-higher remainder after constant and linear terms
are also separated. The full-domain definition supplies
background holomorphy without treating $P_V=1$ as a product
domain.

\subsection{Local coefficient and multiplier bounds}
\label{pm-coefficients-section}

\begin{theorem}[Local measure interaction]\label{pm-measure-theorem}
For the fixed geometry and product background domain above,
the logarithmic interactions
$j^{\rm str}(V;K_Vgz)$, $j^{\rm H}(V;gz)$ and
$v^{\rm G}(V;g,z)$ are holomorphic and have primitive
supports satisfying \eqref{pm-support-bounds}. For
$|g|R<\rho_L$, their sum, regrouped by supports $Y$, obeys
\begin{equation}\label{pm-combined-bound}
 \sup_x\sum_{Y\ni x}e^{a_0|Y|+\kappa\ell(Y)}
 \sup|v_Y(V;g,z)|
 \le C_{L,a_0,\kappa}
   (|g|R+|g|^2R^2+|g|R^3).
\end{equation}
The constants are independent of $M$ and periodic volume.
If the right side is at most one, the same bound multiplied
by $e$ holds with $|v_Y|$ replaced by $|e^{v_Y}-1|$.
For $R\ge1$ and $|g|R\le1$, it is at most
$C_{L,a_0,\kappa}|g|R^3$.
\end{theorem}
\begin{proof}
Combine \eqref{pm-jacobian-bound},
\eqref{pm-haar-bound} and \eqref{pm-gauge-bound} and group
equal supports. The local coefficient assertion can also
be seen directly from the convergent expansions
\[
 \begin{split}
 j_c^{\rm str}(V;K_Vgz)&=\sum_{n\ge1}g^n s_{c,n}(V;z),\\
 j_b^{\rm H}(V;gz)&=\sum_{n\ge2}g^n h_{b,n}(V;z),\\
 v_r^{\rm G}(V;g,z)&=\sum_{n\ge3}g^{n-2}v_{r,n}(V;z).
 \end{split}
\]
Each coefficient is homogeneous of degree $n$ in $z$,
holomorphic in only the background coordinates of its support,
and has local norm at most $C_0(C_1R)^n$ by Cauchy on a
fixed fluctuation disc. Background Cauchy estimates on the
smaller product domain preserve the same supports. These are
anchored coefficient bounds, rather than an extensive bound
on the full sum. Finally $|e^v-1|\le e|v|$ for $|v|\le1$.
\end{proof}

The support-cardinality factor $K^{|Y|}$ in the overlapping
interaction theorem is absorbed by increasing $a_0$ by
$\log\max(1,K)$. This is a fixed-$L$ statement, not a
claim of uniform constants when $L$ or the cell dimension is
allowed to increase. For a logarithmically growing cutoff
$R(g)$, these measure interactions tend to zero as soon as
$|g|R(g)^3\to0$.

The fixed real cutoff is the original retained window
$\|z(b)\|<R$. It remains unchanged because
$\Phi_V(gz)(b)=gz(b)$ on retained bonds. Additional pivot,
contour or reconstructed-field indicators retain their exact
pullbacks and are not covered by this theorem. CMP116,
p.~12 immediately before (2.3), explicitly uses the
pivot-removed bond set; its (2.14), p.~15, keeps this retained
cutoff fixed during Gaussian deformation. CMP109,
pp.~266--268, (2.5) and (2.9)--(2.13), supplies the gauge
and pivot setup. The unreadable product subscript of
CMP109 (2.9) in the accessible OCR is not presented as a
visually verified formula.

These factors are strictly local in independent current
backgrounds $V$. Substituting $V_*(W)$ does not give them
strict finite support in $W$. The full classical remainder,
the previously integrated interaction, the external factor
\eqref{pm-delta-reduction}, and any extra cutoff indicators
remain distinct obligations.

\section{Localization and the next quantum induction}
\label{ps19-next}

The homogeneous physical Gaussian is now identified in the source
pivot coordinates, on a common small complex neighborhood with
constants independent of the preceding classical depth. Its
covariance has the spectral and spatial estimates required by
the marked overlap expansion of Notebook 18. The retained
cutoff coordinates are fixed, and the local measure/gauge
factors have their explicit analytic bounds.

The strictly background-local version requires a small local
decomposition of the actual precision and of the action
remainder; exponential matrix-kernel decay alone does not
supply it. Notebook 20 verifies these hypotheses directly for
the bare Wilson step. Notebook 21 constructs them uniformly
through the minimized classical actions using a rooted support
fixed point, and closes one integration of that classical
reference action. Composition with the stationary response
and all inherited quantum corrections must still be included
in a complete induction.

The original CMP116 localization and integration results provide
one route once their orbit/reconstruction and deformed-operator
hypotheses are matched. Notebook 21 supplies the direct
fixed-point route on our homogeneous small product domain.
The next steps must extend to bad regions and running
interactions while preserving the actual normalization.

That induction, uniform sourced observable bounds, and the
continuum symmetry and reconstruction arguments remain the
steps from these finite-regulator results to four-dimensional
Yang--Mills existence. They are the continuing objective.

\begin{thebibliography}{9}
\bibitem{prior19}
\emph{Toward a four-dimensional Yang--Mills construction},
research Notebooks 5--18 in this workspace. The ordinary nonlinear
minimum, classical coarse response and flat retained-coordinate
covariance are the principal inputs.
\bibitem{balaban10919}
T.~Ba\l aban,
\emph{Renormalization group approach to lattice gauge field
theories I}, Commun. Math. Phys. \textbf{109} (1987), 249--301,
\href{https://scispace.com/pdf/renormalization-group-approach-to-lattice-gauge-field-3ygai0vxt1.pdf}{original article copy},
equations (0.11)--(0.16), (2.3)--(2.13).
\bibitem{balaban11619}
T.~Ba\l aban,
\emph{Renormalization group approach to lattice gauge field
theories. II. Cluster expansions},
Commun. Math. Phys. \textbf{116} (1988), 1--22,
\href{https://scispace.com/pdf/renormalization-group-approach-to-lattice-gauge-field-4penm4k7hc.pdf}{original article copy}.
\bibitem{source19}
E.~Sussman and collaborators,
\href{https://github.com/ethanwsussman/Yang--Mills-LLM-experiment}{Yang--Mills-LLM-experiment},
deposited raw accounts at commit
\texttt{e392ef45013a581664bab76b90867170b30fe10e}.
\end{thebibliography}
\end{document}
